% ファイル構成
\section{ファイル構成}

\subsection{ディレクトリ構造}

\begin{lstlisting}[basicstyle=\ttfamily\small, frame=single]
Embedded-Module-Architecture/
+-- .devcontainer/
|   +-- Dockerfile              ... LaTeX環境用コンテナ定義
|   +-- devcontainer.json       ... DevContainer設定
+-- doc/
|   +-- .latexmkrc              ... latexmk設定
|   +-- main.tex                ... 仕様書本体
|   +-- sections/               ... 仕様書の各セクション
|       +-- cover.tex
|       +-- revision.tex
|       +-- overview.tex
|       +-- background.tex
|       +-- purpose.tex
|       +-- architecture.tex
|       +-- files.tex
|       +-- techstack.tex
|       +-- api.tex
|       +-- execution.tex
+-- include/                    ... プロジェクト固有ヘッダー
|   +-- SystemData.h            ... SystemData構造体（モジュール間共有データ）
|   +-- ProjectConfig.h         ... Configインスタンス定義（ピンアサイン等）
+-- lib/
|   +-- ModuleCore/             ... コアライブラリ
|   |   +-- IModule.h           ... IModuleインターフェース
|   |   +-- ModuleTimer.h       ... ModuleTimerクラス
|   +-- LedModule/              ... モジュール実装例（出力）
|   |   +-- LedModule.h
|   |   +-- LedModule.cpp
|   +-- TftModule/              ... TFT LCDモジュール（出力+タッチ入力）
|   |   +-- TftModule.h
|   |   +-- TftModule.cpp
|   +-- Mpu6500Module/          ... MPU-6500 IMUセンサーモジュール（入力）
|   |   +-- Mpu6500Module.h
|   |   +-- Mpu6500Module.cpp
|   +-- CameraModule/           ... OV5640カメラモジュール（入力）
|   |   +-- CameraModule.h
|   |   +-- CameraModule.cpp
|   +-- ServoModule/            ... サーボモーター制御モジュール（出力）
|   |   +-- ServoModule.h
|   |   +-- ServoModule.cpp
|   +-- DriveMotorModule/       ... DCモーター単体制御モジュール（出力）
|   |   +-- DriveMotorModule.h
|   |   +-- DriveMotorModule.cpp
|   +-- ChassisModule/          ... 4輪統合足回り制御モジュール（出力）
|       +-- ChassisModule.h
|       +-- ChassisModule.cpp
+-- src/
|   +-- main.cpp                ... エントリーポイント
+-- test/                       ... テストコード
+-- platformio.ini              ... PlatformIO設定
+-- README.md                   ... プロジェクト概要
\end{lstlisting}

\subsection{ファイル一覧と機能対応}

\begin{tabularx}{\textwidth}{l l X}
\toprule
\textbf{ファイル} & \textbf{種別} & \textbf{説明} \\
\midrule
\texttt{lib/ModuleCore/IModule.h} & コア & IModuleインターフェース定義 \\
\texttt{lib/ModuleCore/ModuleTimer.h} & コア & ModuleTimerクラス定義 \\
\texttt{lib/LedModule/LedModule.h/.cpp} & 出力モジュール & GPIO LED制御（リファレンス実装） \\
\texttt{lib/TftModule/TftModule.h/.cpp} & 出力+入力 & 2.8インチSPI TFT LCD表示 + XPT2046タッチ入力 \\
\texttt{lib/Mpu6500Module/Mpu6500Module.h/.cpp} & 入力モジュール & MPU-6500 IMUセンサー（I2C直接アクセス） \\
\texttt{lib/CameraModule/CameraModule.h/.cpp} & 入力モジュール & OV5640カメラ（DVP / esp\_camera） \\
\texttt{lib/ServoModule/ServoModule.h/.cpp} & 出力モジュール & サーボモーター制御（LEDC PWM） \\
\texttt{lib/DriveMotorModule/DriveMotorModule.h/.cpp} & 出力モジュール & DCモーター単体制御（Hブリッジ） \\
\texttt{lib/ChassisModule/ChassisModule.h/.cpp} & 出力モジュール & 4輪統合足回り制御（機能統合系） \\
\texttt{include/SystemData.h} & プロジェクト固有 & 各モジュールのData構造体を集約したSystemData構造体 \\
\texttt{include/ProjectConfig.h} & プロジェクト固有 & モジュールConfigインスタンス（ピンアサイン等） \\
\texttt{src/main.cpp} & プロジェクト固有 & エントリーポイント（モジュール登録・ループ制御） \\
\bottomrule
\end{tabularx}

\subsection{ライブラリ構成方針}

\begin{itemize}
    \item \textbf{lib/ModuleCore/} — プロジェクト非依存のコアコンポーネント（IModule、ModuleTimer）。
    他のプロジェクトにもそのまま流用可能。
    \item \textbf{lib/各モジュール名/} — 個別のハードウェアモジュール実装。
    それぞれが\texttt{IModule}を継承し、モジュール固有のサブ構造体を定義する。
    \texttt{.h}（宣言 + 前方宣言）と\texttt{.cpp}（実装）に分離する。
    \item \textbf{include/} — プロジェクト固有のヘッダー
    （\texttt{ProjectConfig.h}等、プロジェクトで使用するモジュールの組み合わせ・ピンアサインに依存するファイル）。
\end{itemize}
